\documentclass[10pt,a4paper]{article}
\usepackage[a4paper,margin=2.1cm]{geometry}
\usepackage[T1]{fontenc}
\usepackage[utf8]{inputenc}
\usepackage{lmodern}
\usepackage{microtype}
\usepackage{xcolor}
\usepackage{enumitem}
\usepackage{tabularx}
\usepackage{longtable}
\usepackage{array}
\usepackage{amssymb}
\usepackage{hyperref}
\usepackage{fancyhdr}
\definecolor{tudblue}{RGB}{0,90,140}
\definecolor{draftred}{RGB}{165,35,35}
\definecolor{lightblue}{RGB}{232,242,247}
\hypersetup{colorlinks=true,linkcolor=tudblue,urlcolor=tudblue}
\setlength{\parindent}{0pt}
\setlength{\parskip}{0.55em}
\setlist{itemsep=0.25em,topsep=0.35em}
\renewcommand{\arraystretch}{1.22}
\pagestyle{fancy}
\fancyhf{}
\lhead{TU Darmstadt Ethics Commission}
\rhead{Draft for supervisor review}
\cfoot{\thepage}
\newcommand{\checked}{\ensuremath{\boxtimes}}
\newcommand{\unchecked}{\ensuremath{\square}}
\newcommand{\confirm}[1]{\textcolor{draftred}{\textbf{TO CONFIRM:} #1}}
\newcommand{\field}[2]{\textbf{#1}\par #2\par}
\newcommand{\draftbanner}{%
  \begin{center}
  \fcolorbox{draftred}{white}{\parbox{0.91\textwidth}{\centering
  \textcolor{draftred}{\textbf{DRAFT -- NOT FOR SUBMISSION OR DATA COLLECTION}}\\
  Administrative and data-management fields marked ``TO CONFIRM'' require supervisor approval.}}
  \end{center}}
\newcommand{\sectionbox}[1]{%
  \vspace{0.4em}\noindent\colorbox{lightblue}{\parbox{\dimexpr\linewidth-2\fboxsep\relax}{\textbf{#1}}}\vspace{0.35em}}

\newcommand{\scale}[1]{\textit{Response scale: #1}}
\newcommand{\qitem}[2]{\textbf{#1} #2\par}
\begin{document}

\begin{center}
  {\Large\bfseries Questionnaire and Study Measures}\\[0.3em]
  {\large Local Trait Inference on Smartphones}
\end{center}
\draftbanner

This appendix lists all participant-facing questionnaire items planned for the Android study. Scale values are stored numerically together with their verbal labels. The optional camera module is a data-collection instrument, not a questionnaire.

\section*{A. Eligibility and Demographic Questions}

\qitem{A1. Age confirmation}{I confirm that I am at least 18 years old.}
\textit{Options: Yes / No. A ``No'' response ends eligibility.}

\qitem{A2. Age group}{Which age group are you in?}
\textit{Options: 18--24; 25--34; 35--44; 45--54; 55 or older; prefer not to say.}

\qitem{A3. Gender}{How do you describe your gender?}
\textit{Options: woman; man; non-binary; self-description; prefer not to say.}

\qitem{A4. Highest completed education level}{What is the highest level of education you have completed?}
\textit{Options: school qualification; vocational training; bachelor's degree; master's degree or higher; other; prefer not to say.}

\textbf{Implementation alignment note:} The current prototype lacks ``prefer not to say'' for age, and lacks ``self-description''/``other'' options in two items. These options should be aligned before data collection.

\section*{B. Internet Users' Information Privacy Concerns -- IUIPC-8}

Instruction: The following statements concern online privacy. Please indicate how strongly you agree or disagree with each statement.

\scale{1 strongly disagree; 2 disagree; 3 somewhat disagree; 4 neither agree nor disagree; 5 somewhat agree; 6 agree; 7 strongly agree.}

\subsection*{Control}
\qitem{B1.}{Consumer online privacy is really a matter of consumers' right to exercise control and autonomy over decisions about how their information is collected, used, and shared.}
\qitem{B2.}{Consumer control of personal information lies at the heart of consumer privacy.}

\subsection*{Awareness}
\qitem{B3.}{Companies seeking information online should disclose the way the data are collected, processed, and used.}
\qitem{B4.}{A good consumer online privacy policy should have a clear and conspicuous disclosure.}

\subsection*{Collection}
\qitem{B5.}{It usually bothers me when online companies ask me for personal information.}
\qitem{B6.}{When online companies ask me for personal information, I sometimes think twice before providing it.}
\qitem{B7.}{It bothers me to give personal information to so many online companies.}
\qitem{B8.}{I am concerned that online companies are collecting too much personal information about me.}

\textbf{Scoring note for researchers:} Preserve the three dimensions Control, Awareness, and Collection. Confirm the planned scoring and psychometric analysis with the psychology supervisor before preregistration or final analysis.

\textbf{Source:} Gross, T. (2023). Toward Valid and Reliable Privacy Concern Scales: The Example of IUIPC-8. In \textit{Human Factors in Privacy Research}. The eight-item form omits ctrl3 and awa3 from IUIPC-10 and uses a seven-point agreement scale.

\newpage
\section*{C. Affinity for Technology Interaction -- ATI}

Instruction: In the following questionnaire, we will ask you about your interaction with technical systems. The term ``technical systems'' refers to apps and other software applications, as well as entire digital devices such as a mobile phone, computer, TV, or car navigation system.

\scale{1 completely disagree; 2 largely disagree; 3 slightly disagree; 4 slightly agree; 5 largely agree; 6 completely agree.}

\qitem{C1.}{I like to occupy myself in greater detail with technical systems.}
\qitem{C2.}{I like testing the functions of new technical systems.}
\qitem{C3.}{I predominantly deal with technical systems because I have to.}
\qitem{C4.}{When I have a new technical system in front of me, I try it out intensively.}
\qitem{C5.}{I enjoy spending time becoming acquainted with a new technical system.}
\qitem{C6.}{It is enough for me that a technical system works; I do not care how or why.}
\qitem{C7.}{I try to understand how a technical system exactly works.}
\qitem{C8.}{It is enough for me to know the basic functions of a technical system.}
\qitem{C9.}{I try to make full use of the capabilities of a technical system.}

\textbf{Scoring note for researchers:} Items C3, C6, and C8 are reverse-coded. Compute the mean across all nine items after reversal. Report the mean, standard deviation, and internal consistency as appropriate.

\textbf{Source:} Franke, T., Attig, C., \& Wessel, D. (2019). A Personal Resource for Technology Interaction: Development and Validation of the Affinity for Technology Interaction (ATI) Scale. \textit{International Journal of Human--Computer Interaction, 35}(6), 456--467. \url{https://doi.org/10.1080/10447318.2018.1456150}

\section*{D. Baseline Local-AI Expectations}

Instruction: Please indicate how strongly you agree or disagree before seeing any file analysis.

\scale{1 strongly disagree; 2 disagree; 3 neither agree nor disagree; 4 agree; 5 strongly agree.}

\qitem{D1.}{I expect a phone to infer sensitive information from selected local files.}
\qitem{D2.}{The statement ``data stays on the phone'' makes this feel privacy-friendly.}
\qitem{D3.}{I usually enjoy trying new digital technologies.}
\qitem{D4.}{I am concerned about how apps can use information stored on my phone.}

\textbf{Design note:} D3 overlaps conceptually with ATI and D4 overlaps with IUIPC-8. Confirm whether they are needed as brief face-valid items or should be removed to reduce burden and avoid redundant measurement.

\section*{E. Post-Inference Questions}

These three questions are shown after each successfully analyzed image or PDF first-page preview. The displayed probabilistic inference remains visible while the participant responds.

\scale{1 not at all; 2 slightly; 3 moderately; 4 strongly; 5 very strongly.}

\qitem{E1.}{How unexpected was this inference?}
\qitem{E2.}{How sensitive does this inference feel?}
\qitem{E3.}{How accurate does this inference seem?}

Each response is linked to a randomized file identifier and its model output. Original filenames and local paths are excluded from the final research export.

\newpage
\section*{F. Combined Profile and Final Questions}

Before these questions, participants see one compact probabilistic personal profile synthesized from the file-level inferences. The following framing text is displayed:

\begin{quote}
Imagine a normal app performed this kind of local analysis regularly in the background and only sent inferred summaries or profile categories away.
\end{quote}

\scale{1 strongly disagree; 2 disagree; 3 neither agree nor disagree; 4 agree; 5 strongly agree.}

\qitem{F1.}{Before this study, I expected my phone to be able to run this kind of AI analysis locally.}
\qitem{F2.}{I would be concerned if an app repeated this analysis regularly on local photos, screenshots, or documents.}
\qitem{F3.}{The statement ``the original files never leave the device'' would still reassure me if inferred summaries or profile categories were sent to a server.}
\qitem{F4.}{I would be comfortable if an app analyzed local files in the background and sent inferred summaries to a server without clearly notifying me each time.}
\qitem{F5.}{I would be comfortable if this inferred profile were used commercially, for example for advertisements, product offers, or sharing with another company.}
\qitem{F6.}{After seeing this profile, I would still allow broad photo or file access to an app.}

\section*{G. Optional Front-Camera Reaction Measure}

This is not a self-report questionnaire. After separate opt-in, the app may sample a low-resolution front-camera frame approximately every 12 seconds while the relevant study interface is active. Local classification yields:
\begin{itemize}
  \item timestamp;
  \item study-screen event;
  \item model identifier/version;
  \item predicted label from angry, disgust, fear, happy, neutral, sad, or surprise; and
  \item model confidence.
\end{itemize}

No audio is captured. The final implementation must delete each temporary frame immediately and retain no image bytes or file path. This measure is exploratory and must not be interpreted as a direct measure of the participant's true emotional state or as a diagnosis.

\section*{H. End-of-Study Information}

Participants are reminded that:
\begin{itemize}
  \item all displayed traits are probabilistic model outputs and may be wrong;
  \item the original selected files were processed locally and were not part of the research export;
  \item local processing does not by itself prevent privacy impact if derived summaries or profiles are transmitted; and
  \item they may contact the study team with questions or request deletion using their study ID until the stated cutoff.
\end{itemize}

\end{document}
